\chapter*{Реферат}
\addcontentsline{toc}{chapter}{Реферат}
{\centering\bfseries\WorkingTitle\par}
\medskip
{\centering\AuthorName\par}
\medskip
\noindent\textbf{Ключевые слова:} вызванный составной потенциал действия (ECAP), стимуляция спинного мозга, коннектом, \textit{Drosophila melanogaster}, нейродинамика, предобучение, перенос представлений.
\medskip
Работа посвящена биофизическому моделированию кривой роста ECAP. Цель --- разработать метод с предобучением временного модуля на коннектомной динамике \textit{Drosophila melanogaster} и оценить вклад предобучения в точность моделирования человеческого электрического ответа.

В первой главе сопоставлены наблюдаемые реакции Drosophila, ограничения сетевых моделей, механизмы формирования ECAP и методы предобучения. Определены требования к совместимости исходных наблюдений и независимой проверке. Экспериментальные результаты предлагаемого переноса не представлены.

Планируется разработать генератор нейродинамических рядов, определить переносимую часть параметров и сопрячь временной модуль с отдельным человеческим биофизическим оператором. Предусматриваются проверка оператора и сравнение с обучением без переноса и простой кривой роста. Будет разграничено влияние структуры, динамики и предобучения. Основная единица оценки --- участник, показатель --- ошибка амплитуды кривой роста.

Применение ограничено численным исследованием измерительного канала. ECAP и нейродинамика Drosophila не являются прямыми измерениями субъективной боли; клинические выводы потребуют отдельной проверки.
